% Generated balanced MCMC validation summary table.
\begin{table}[!htbp]
\centering
\scriptsize
\resizebox{\textwidth}{!}{%
\begin{tabular}{@{}>{\raggedright\arraybackslash}p{0.20\textwidth}>{\raggedright\arraybackslash}p{0.30\textwidth}rrrr@{}}
\toprule
Scenario & Mechanism & \shortstack{Joint QDESN\\\(\AL\)--\RHS} & \shortstack{Independent QDESN\\\(\AL\)--\RHS} & \shortstack{Joint exQDESN\\\(\exAL\)--\RHS} & \shortstack{Independent exQDESN\\\(\exAL\)--\RHS}
\\
\midrule
Asymmetric Laplace Tail & Asymmetric Laplace; seasonal AR(1) location-scale & 0.099 (0.084; 1) & \textbf{0.083 (0.088; 7)} & 0.093 (0.081; 0) & 0.089 (0.077; 0) \\
Gaussian-Mixture Bridge & Gaussian mixture; seasonal AR(1) location-scale & 0.096 (0.097; 0) & 0.099 (0.101; 0) & 0.106 (0.096; 0) & \textbf{0.094 (0.095; 0)} \\
Laplace Bridge & Laplace; seasonal AR(1) location-scale & \textbf{0.085 (0.089; 0)} & 0.094 (0.091; 7) & 0.111 (0.113; 0) & 0.106 (0.117; 0) \\
Nonlinear Reservoir Friendly & Gaussian mixture; nonlinear reservoir dynamics & \textbf{0.104 (0.109; 0)} & 0.116 (0.112; 1) & 0.157 (0.151; 0) & 0.149 (0.143; 0) \\
Normal Bridge & Gaussian; seasonal AR(1) location-scale & \textbf{0.081 (0.082; 0)} & 0.082 (0.082; 2) & 0.115 (0.097; 0) & 0.104 (0.101; 0) \\
Persistent Heavy Tail & Student-t; persistent seasonal AR(1) dynamics & 0.131 (0.120; 0) & 0.124 (0.118; 0) & 0.130 (0.103; 0) & \textbf{0.106 (0.100; 0)} \\
Regime Shift & Student-t; regime-shift location-scale & 0.153 (0.085; 0) & \textbf{0.096 (0.091; 6)} & 0.185 (0.090; 0) & 0.113 (0.087; 0) \\
Student-t Location-Scale & Student-t; seasonal AR(1) location-scale & 0.080 (0.062; 0) & \textbf{0.071 (0.065; 1)} & 0.116 (0.093; 0) & 0.099 (0.092; 0) \\
\bottomrule
\end{tabular}
}%
\caption{Scenario-level MCMC validation analysis for the joint multi-quantile validation study. Entries are forecast MAE with fit-window MAE and the number of raw forecast-window adjacent-level crossings in parentheses, written as forecast MAE (fit MAE; raw crossings). Boldface marks the lowest forecast MAE within each pre-specified eight-mechanism scenario. QDESN uses the \(\AL\) working likelihood, exQDESN uses the \(\exAL\) working likelihood, and \RHS{} denotes the regularized horseshoe prior. All scores use the monotone quantile-grid reporting rule; raw crossings are retained only as pre-rearrangement diagnostics. exAL rows are included when finite-path and quantile-stability criteria are satisfied under the balanced validation design.}
\label{tab:joint-qdesn-article-validation-mcmc-balanced-model-summary}
\end{table}
